\documentclass[11pt,a4paper]{article}

\usepackage[utf8]{inputenc}
\usepackage[T2A]{fontenc}
\usepackage[ukrainian,english]{babel}
\usepackage[top=22mm,bottom=22mm,left=22mm,right=22mm]{geometry}
\usepackage{hyperref}
\usepackage{fancyhdr}
\usepackage{parskip}
\usepackage{tcolorbox}
\usepackage{enumitem}
\usepackage{listings}
\usepackage{xcolor}

\tcbuselibrary{skins,breakable}

\hypersetup{
    pdftitle={Workshop 08 --- Headless / CI/CD: Handout},
    pdfauthor={Vadym Bondarenko},
    colorlinks=true,
    linkcolor=blue,
    urlcolor=cyan,
}

\pagestyle{fancy}
\fancyhf{}
\fancyhead[L]{\small\textit{Workshop 08 --- Headless / CI/CD}}
\fancyhead[R]{\small\thepage}
\fancyfoot[C]{\small\textit{vcdev.me \textperiodcentered{} 2026}}
\renewcommand{\headrulewidth}{0.4pt}

\lstdefinelanguage{yaml}{
    keywords={true,false,null},
    sensitive=false,
    comment=[l]{\#},
    morestring=[b]',
    morestring=[b]"
}

\lstset{
    basicstyle=\ttfamily\small,
    breaklines=true,
    frame=single,
    backgroundcolor=\color{gray!10},
    rulecolor=\color{gray!40},
    xleftmargin=8pt,
    xrightmargin=8pt,
    aboveskip=6pt,
    belowskip=6pt,
}

\title{%
  {\Huge\bfseries Workshop 08: Headless / CI/CD}\\[0.6em]
  {\large Productionizing \texttt{claude -p} --- handout}
}
\author{Vadym Bondarenko \\ \small\href{mailto:vadym.bondarenko@vcdev.me}{vadym.bondarenko@vcdev.me}}
\date{2026}

\begin{document}

\maketitle

\begin{abstract}
Пост-воркшоп референс. Не повторює слайди --- сухе зведення CLI-flag-ів, output форматів, permission-режимів, auth-precedence, CI templates і чек-листів. Усі факти посилаються на офіційні docs (\href{https://code.claude.com/docs/en/headless}{code.claude.com}).
\end{abstract}

\tableofcontents
\newpage

\section{Headless mode: основа}

\texttt{claude -p "<prompt>"} (long form \texttt{--print}) --- non-interactive mode, exit після відповіді. Раніше називалось ``headless''; тепер просто \texttt{-p}.

\begin{lstlisting}
claude -p "What does the auth module do?"
claude -p "Summarize this project" --output-format json | jq -r '.result'
claude -p "Fix tests" --max-turns 5 --allowedTools "Bash,Read,Edit"
\end{lstlisting}

\textbf{Що НЕ працює у \texttt{-p}:}
\begin{itemize}[leftmargin=2em]
\item Permission prompts --- нема UI; pre-approve через \texttt{--allowedTools} або \texttt{--permission-mode}
\item Skills типу \texttt{/commit} --- не доступні у \texttt{-p}; описуй прозою
\item Auto mode --- після 3 блоків поспіль аборт
\end{itemize}

\section{\texttt{--bare} mode}

\texttt{claude --bare -p "..."} пропускає auto-discovery: hooks, skills, plugins, MCP servers, auto memory, \texttt{CLAUDE.md}.

\begin{itemize}[leftmargin=2em]
\item Однаковий результат на будь-якій машині (важливо для CI)
\item Інструменти: Bash, Read, Edit (інші додавай явно)
\item \textbf{Caveat:} bare mode НЕ читає \texttt{CLAUDE\_CODE\_OAUTH\_TOKEN}. Лише \texttt{ANTHROPIC\_API\_KEY} або \texttt{apiKeyHelper}
\end{itemize}

\section{Output formats}

\subsection{\texttt{--output-format text} (default)}

Plain text у stdout.

\subsection{\texttt{--output-format json}}

\begin{lstlisting}
claude -p "Summarize" --output-format json | jq -r '.result'
\end{lstlisting}

Поля JSON-відповіді:
\begin{itemize}[leftmargin=2em]
\item \texttt{result} --- текстова відповідь Claude
\item \texttt{session\_id} --- UUID, для \texttt{--resume}
\item \texttt{total\_cost\_usd} --- витрати на запит (моніторинг!)
\item \texttt{num\_turns} --- скільки ітерацій
\item \texttt{is\_error} --- bool
\item \texttt{structured\_output} --- якщо передано \texttt{--json-schema}
\end{itemize}

\subsection{\texttt{--output-format stream-json}}

Newline-delimited events. Типи з docs:
\begin{itemize}[leftmargin=2em]
\item \texttt{system/init} --- перша подія; \texttt{plugins[]} і \texttt{plugin\_errors[]} (fail CI якщо плагін не завантажився)
\item \texttt{system/api\_retry} --- перед ретраєм; \texttt{attempt}, \texttt{max\_retries}, \texttt{retry\_delay\_ms}, \texttt{error\_status}, \texttt{error}
\item \texttt{stream\_event} --- з \texttt{--include-partial-messages}; partial deltas
\end{itemize}

\subsection{\texttt{--json-schema}}

Гарантія структури:

\begin{lstlisting}
claude -p "Extract function names from auth.py" \
  --output-format json \
  --json-schema '{"type":"object","properties":{
    "functions":{"type":"array","items":{"type":"string"}}},
    "required":["functions"]}' \
  | jq '.structured_output'
\end{lstlisting}

\section{Permissions у headless}

\subsection{6 permission-режимів}

\begin{tabular}{lll}
\hline
\textbf{Mode} & \textbf{Що автоприймає} & \textbf{CI use case} \\
\hline
\texttt{default} & Reads only & Sensitive review \\
\texttt{acceptEdits} & Reads + edits + filesystem cmds & Apply-fix у MR \\
\texttt{plan} & Reads only & Readonly research \\
\texttt{auto} & Майже все, з classifier & \textbf{Не для \texttt{-p}} \\
\texttt{dontAsk} & Тільки allow-rules & \textbf{Locked-down CI} \\
\texttt{bypassPermissions} & Все, окрім protected paths & Container/VM only \\
\hline
\end{tabular}

\subsection{Permission rule syntax}

\texttt{Tool(matcher)} --- patterns використовують glob \texttt{*}.

\textbf{Critical:} пробіл перед \texttt{*} обовʼязковий.

\begin{lstlisting}
# WRONG: also matches `git diff-index`
--allowedTools "Bash(git diff*)"

# RIGHT: only matches `git diff <something>`
--allowedTools "Bash(git diff *)"
\end{lstlisting}

\subsection{settings.json для CI}

\begin{lstlisting}
{
  "permissions": {
    "allow": [
      "Bash(git *)",
      "Read(.)"
    ],
    "deny": [
      "Read(./.env)",
      "Read(./.env.*)",
      "Read(./secrets/**)",
      "WebFetch",
      "Bash(curl *)"
    ],
    "defaultMode": "dontAsk"
  }
}
\end{lstlisting}

\textbf{Запуск:} \texttt{claude --settings ./ci-settings.json -p "..."}.

\subsection{Protected paths}

Завжди заборонені для авто-запису: \texttt{.git}, \texttt{.vscode}, \texttt{.idea}, \texttt{.husky}, \texttt{.claude} (з винятками для \texttt{commands/}, \texttt{agents/}, \texttt{skills/}, \texttt{worktrees/}), \texttt{.gitconfig}, \texttt{.bashrc}/\texttt{.zshrc}, \texttt{.mcp.json}, \texttt{.claude.json}.

\section{Auth у CI}

\subsection{Precedence (від високого до низького)}

\begin{enumerate}
\item Cloud (\texttt{CLAUDE\_CODE\_USE\_BEDROCK} / \texttt{\_VERTEX} / \texttt{\_FOUNDRY})
\item \texttt{ANTHROPIC\_AUTH\_TOKEN} --- Bearer (для LLM-gateway proxy)
\item \texttt{ANTHROPIC\_API\_KEY} --- X-Api-Key (direct Anthropic API)
\item \texttt{apiKeyHelper} --- script для rotating creds
\item \texttt{CLAUDE\_CODE\_OAUTH\_TOKEN} --- long-lived OAuth (subscription)
\item \texttt{/login} OAuth --- default Pro/Max/Team/Enterprise
\end{enumerate}

\subsection{\texttt{claude setup-token}}

Для CI без браузера. 1 рік валідності, subscription auth.

\begin{lstlisting}
claude setup-token
# walks through OAuth, prints token to terminal
export CLAUDE_CODE_OAUTH_TOKEN=<token>
\end{lstlisting}

\textbf{Caveat:} \texttt{--bare} НЕ читає цю змінну.

\subsection{\texttt{apiKeyHelper}}

\begin{lstlisting}
{ "apiKeyHelper": "/usr/local/bin/get-api-key.sh" }
\end{lstlisting}

Скрипт виводить ключ у stdout. Викликається кожні 5 хв або на HTTP 401. Тюнинг: \texttt{CLAUDE\_CODE\_API\_KEY\_HELPER\_TTL\_MS}.

\subsection{OIDC для cloud}

\begin{itemize}[leftmargin=2em]
\item GitLab \texttt{->} AWS: \texttt{aws sts assume-role-with-web-identity} через \texttt{CI\_JOB\_JWT\_V2}
\item GitLab \texttt{->} GCP: WIF + \texttt{gcloud auth login --cred-file=<(...)}
\item GitHub \texttt{->} AWS: \texttt{aws-actions/configure-aws-credentials@v4} + \texttt{id-token: write}
\item GitHub \texttt{->} GCP: \texttt{google-github-actions/auth@v2}
\end{itemize}

\section{Cost ceilings}

\begin{tabular}{ll}
\hline
\textbf{Flag} & \textbf{Семантика} \\
\hline
\texttt{--max-turns N} & N agentic ітерацій \texttt{->} exit з error. Default: no limit \\
\texttt{--max-budget-usd \$} & Hard stop по витратах \\
\texttt{--fallback-model} & Fallback на secondary якщо primary overloaded \\
\texttt{--no-session-persistence} & Не пише сесію на диск \\
\hline
\end{tabular}

\section{GitLab CI шаблон}

\begin{lstlisting}[language=yaml]
stages:
  - ai

claude-review:
  stage: ai
  image: node:24-alpine3.21
  timeout: 10 minutes
  rules:
    - if: '$CI_PIPELINE_SOURCE == "merge_request_event"'
  before_script:
    - apk add --no-cache git curl bash jq
    - curl -fsSL https://claude.ai/install.sh | bash
    - export PATH="$HOME/.local/bin:$PATH"
  script:
    - git fetch origin "$CI_MERGE_REQUEST_TARGET_BRANCH_NAME"
    - git diff "origin/$CI_MERGE_REQUEST_TARGET_BRANCH_NAME"...HEAD > diff.patch
    - >
      cat diff.patch | claude -p
      --bare
      --permission-mode plan
      --output-format json
      --max-turns 3
      --max-budget-usd 1.00
      --append-system-prompt "Senior code reviewer. JSON output."
      > review.json
    - jq -r '.result' review.json
\end{lstlisting}

\textbf{CI/CD variables:} \texttt{ANTHROPIC\_API\_KEY} (masked); опційно \texttt{GITLAB\_ACCESS\_TOKEN} для постингу коментарів.

\section{GitHub Actions шаблон}

\begin{lstlisting}[language=yaml]
name: Claude Review
on:
  pull_request:
    types: [opened, synchronize]
permissions:
  contents: read
  pull-requests: write
jobs:
  review:
    runs-on: ubuntu-latest
    steps:
      - uses: actions/checkout@v4
      - uses: anthropics/claude-code-action@v1
        with:
          anthropic_api_key: ${{ secrets.ANTHROPIC_API_KEY }}
          prompt: "Review for bugs and security."
          claude_args: "--max-turns 5 --max-budget-usd 1.00"
\end{lstlisting}

Action \texttt{anthropics/claude-code-action@v1} приймає:
\begin{itemize}[leftmargin=2em]
\item \texttt{prompt} --- інструкції
\item \texttt{claude\_args} --- pass-through для CLI flag-ів
\item \texttt{anthropic\_api\_key} --- required для direct API
\item \texttt{use\_bedrock} / \texttt{use\_vertex} --- boolean
\end{itemize}

\section{Pitfalls (top-10)}

\begin{enumerate}
\item Interactive auth attempt --- нема ключа \texttt{->} browser \texttt{->} CI hangs
\item Secrets у логах через \texttt{--debug}
\item Permission prompt без allowlist \texttt{->} hang до timeout
\item Runaway costs без \texttt{--max-turns}
\item OAuth token не працює з \texttt{--bare}
\item Auto mode у \texttt{-p} \texttt{->} aborts
\item \texttt{ANTHROPIC\_API\_KEY} переб'є subscription
\item Plugin failures silent без stream-json parsing
\item Skills типу \texttt{/commit} не доступні у \texttt{-p}
\item \texttt{--bare} пропускає \texttt{CLAUDE.md} \texttt{->} project conventions not loaded
\end{enumerate}

\section{Безпечні дефолти для CI}

\begin{lstlisting}
claude -p "$PROMPT" \
  --bare \
  --permission-mode dontAsk \
  --allowedTools "Read,Bash(git diff *),Bash(git log *)" \
  --max-turns 5 \
  --max-budget-usd 1.00 \
  --output-format json \
  --no-session-persistence
\end{lstlisting}

\section{Production чек-ліст}

\begin{itemize}[leftmargin=2em]
\item[$\square$] Auth не в коді: \texttt{secrets} / masked variables
\item[$\square$] \texttt{--max-turns N} ставить стелю
\item[$\square$] \texttt{--max-budget-usd} для hard stop
\item[$\square$] \texttt{--permission-mode dontAsk} або \texttt{acceptEdits}
\item[$\square$] \texttt{--allowedTools} narrow --- нема \texttt{Bash(*)}
\item[$\square$] Deny \texttt{.env}, \texttt{secrets/**} у settings
\item[$\square$] \texttt{--bare} для reproducibility
\item[$\square$] \texttt{--output-format json} для downstream
\item[$\square$] Stream-json \texttt{system/init} parsing якщо плагіни критичні
\item[$\square$] Job timeout на CI-рівні
\item[$\square$] Concurrency limit
\item[$\square$] OIDC замість long-lived keys для cloud
\end{itemize}

\section{Дебаг checklist}

\begin{enumerate}
\item \texttt{claude --version} --- встановлений у CI image?
\item \texttt{echo \$\{ANTHROPIC\_API\_KEY:0:8\}...} --- env var присутній (без друку повного ключа)
\item \texttt{claude auth status} --- exit 0 = OK, JSON-вивід
\item \texttt{--verbose} локально (НЕ в CI без masking)
\item Stream-json \texttt{system/init} --- перевір \texttt{plugin\_errors}
\item \texttt{--max-turns 1} --- мінімізуй для повторюваності
\end{enumerate}

\section{Команди для вправ}

\subsection{Вправа 1 --- claude -p locally}

\begin{lstlisting}
claude -p "Summarize this file" sample.txt --output-format json | jq -r '.result'
claude --bare -p "Summarize" sample.txt --allowedTools "Read"
\end{lstlisting}

\subsection{Вправа 2 --- diff review з JSON Schema}

\begin{lstlisting}
cat diff.patch | claude -p \
  --append-system-prompt "Senior reviewer. Output JSON." \
  --permission-mode plan \
  --output-format json \
  --json-schema "$(cat schema.json)" \
  --max-turns 1 \
  > review.json
\end{lstlisting}

\subsection{Вправа 3 --- GitLab CI}

\texttt{ANTHROPIC\_API\_KEY} (masked, protected); copy \texttt{gitlab-ci-starter.yml} \texttt{->} \texttt{.gitlab-ci.yml}; push до feature branch; open MR.

\subsection{Вправа 4 --- mock OIDC}

\begin{lstlisting}
bash mock-aws-oidc.sh   # симулятор aws sts assume-role-with-web-identity
\end{lstlisting}

\section{Resources}

\begin{itemize}
\item \href{https://code.claude.com/docs/en/headless}{code.claude.com/docs/en/headless} --- \texttt{-p} flag, \texttt{--bare}, output formats
\item \href{https://code.claude.com/docs/en/cli-reference}{code.claude.com/docs/en/cli-reference} --- повний список flag-ів
\item \href{https://code.claude.com/docs/en/permission-modes}{code.claude.com/docs/en/permission-modes} --- 6 режимів, classifier behavior
\item \href{https://code.claude.com/docs/en/authentication}{code.claude.com/docs/en/authentication} --- precedence, \texttt{setup-token}, \texttt{apiKeyHelper}
\item \href{https://code.claude.com/docs/en/settings}{code.claude.com/docs/en/settings} --- settings.json schema
\item \href{https://code.claude.com/docs/en/gitlab-ci-cd}{code.claude.com/docs/en/gitlab-ci-cd} --- GitLab templates
\item \href{https://code.claude.com/docs/en/github-actions}{code.claude.com/docs/en/github-actions} --- GitHub Actions v1
\item \href{https://github.com/anthropics/claude-code-action}{github.com/anthropics/claude-code-action} --- офіційний action
\item Цей репо: \texttt{workshops/08-headless-ci/exercises/} і \texttt{solutions/}
\end{itemize}

\end{document}
